\documentclass[11pt]{article}

\usepackage[margin=1in]{geometry}
\usepackage{amsmath, amssymb}
\usepackage{booktabs}
\usepackage{enumitem}

\setlength{\parindent}{0pt}
\setlength{\parskip}{0.8em}

\begin{document}

\begin{center}
{\Large \textbf{Spot rate \& forward rate practice HW}}\\
\vspace{0.3em}
\end{center}

\begin{enumerate}[label=\textbf{\arabic*.}]

%-----------------------------------------

\item You are given the following term structure of interest rates.

\begin{center}
\begin{tabular}{cc}
\toprule
Years & Annual Effective Spot Rate \\
\midrule
1 & 3.0\% \\
2 & 3.5\% \\
3 & 3.9\% \\
4 & 4.4\% \\
5 & 5.2\% \\
\bottomrule
\end{tabular}
\end{center}

In one year, a 3-year annuity-immediate will be issued. The annuity makes annual payments of 10,000.

Calculate the present value of this annuity one year from now.

\begin{itemize}
\item A. 25,847
\item B. 26,669
\item C. 27,469
\item D. 27,960
\item E. 28,798
\end{itemize}

\medskip
\hrule
\medskip

%-----------------------------------------

\item An investment has the following payment structure:

\[
1{,}000 \text{ payable in one year, } 1{,}000 \text{ payable in two years, and } 1{,}000 \text{ payable in three years.}
\]

The term structure of interest rates is:

\begin{center}
\begin{tabular}{cc}
\toprule
Length of Investment & Interest Rate \\
\midrule
1 year & 6\% \\
2 year & 7\% \\
3 year & 8\% \\
\bottomrule
\end{tabular}
\end{center}

If the investment is purchased now at its value based on the term structure of interest rates, what is the overall yield rate on this investment?

\begin{itemize}
\item A. 7.1\%
\item B. 7.2\%
\item C. 7.3\%
\item D. 7.4\%
\item E. 7.5\%
\end{itemize}


\medskip
\hrule
\medskip



%-----------------------------------------

\item The one-year forward rates, deferred $t$ years, are estimated to be:

\begin{center}
\begin{tabular}{cccccc}
\toprule
Year ($t$) & 0 & 1 & 2 & 3 & 4 \\
\midrule
Forward Rate & 4\% & 6\% & 8\% & 10\% & 12\% \\
\bottomrule
\end{tabular}
\end{center}

Calculate the spot rate for a zero-coupon bond maturing three years from now.

\begin{itemize}
\item A. 4\%
\item B. 5\%
\item C. 6\%
\item D. 7\%
\item E. 8\%
\end{itemize}

\medskip
\hrule
\medskip

%-----------------------------------------

\item A bond matures for its par value of 1,000 in 3 years. It pays annual coupons of 5\%.

The one-, two-, and three-year annual effective spot rates are $X$, $X + 0.01$, and $X + 0.02$, respectively.

The annual effective yield of the bond is 6.93\%.

Calculate $X$.

\begin{itemize}
\item A. 0.045
\item B. 0.050
\item C. 0.055
\item D. 0.060
\item E. 0.065
\end{itemize}

\medskip
\hrule
\medskip

%-----------------------------------------

\item The following term structure of interest rates is based on available bond market data as of 01/01/2023:

\begin{center}
\begin{tabular}{cc}
\toprule
Term to maturity (years) & Spot Rate \\
\midrule
1 & 3.0\% \\
2 & 3.5\% \\
3 & 4.0\% \\
\bottomrule
\end{tabular}
\end{center}

A three-year 1000 face amount bond with an annual coupon rate of 10\% paid annually is issued on 01/01/2023.

Calculate the yield to maturity for this bond.

\begin{itemize}
\item A. 3.82\%
\item B. 3.86\%
\item C. 3.90\%
\item D. 3.94\%
\item E. 3.98\%
\end{itemize}

\end{enumerate}

\end{document}
